\scanpage{207}

ou aux couples $(\alpha,\beta)$ d'éléments de $\widehat{\Pi}_{0,3}$
satisfaisant la relation des lacets (37)
\[
 (97)=(37)\qquad
 \unsure{(\alpha\rho(\alpha)^{-1})}
 =(\beta\sigma_\infty(\beta)^{-1})\ell_1^{\unsure{\nu}}
\]
(qu'on complètera alors en prenant $f=\beta$).

\margin{$\alpha,\beta$ \unsure{déterminés} \unsure{mod}
$(\alpha,\beta)\longmapsto
(\ell_0^{\unsure{n}}\alpha,\ell_0^{\unsure{n}}\beta)$.}

Au changement près de $\alpha$ en $\ell_0^\nu\alpha=\underline{\alpha}_0$,
c'étaient les conventions du \unsure{§ 44|§ 4 h} (où la formule prend
\unsure{donc} la forme
\[
 \underbrace{\underline{\alpha}_0\rho(\underline{\alpha}_0)^{-1}}_{\xi}
 =\ell_0^\nu
 \underbrace{(f\sigma_\infty(f)^{-1})}_{\unsure{\varphi=\eta}}
\]
\unsure{).}

Compte tenu de (79), les conditions

\deleted{\uline{2° Cas}
\unsure{(98)} $g=1$ i.e. $\unsure{\alpha_0}=1$ i.e. $f=\alpha$,
i.e. $\xi=f\rho(f)^{-1}$ i.e.}

(93) (94) prennent aussi la forme
\[
 (98)\qquad u_{\alpha,\beta,f}
 =\mu(t\sigma_\infty+u)\tau_\infty^{\alpha_2},
\]
et impliquent que $\underline{u}=u_{\alpha,\beta,f}$
\deleted{\illegible{1}} normalise la sous-algèbre de l'algèbre des matrices
$M_2(\widehat{\mathbb Z})$ \uline{engendrée} par $\sigma_\infty$
(compte tenu de la relation
\[
 (99)\qquad \tau_\infty(\sigma_\infty)=-\sigma_\infty
 \qquad\text{dans }\mathrm{Gl}(2,\widehat{\mathbb Z}).
\]

Si on regarde le sous-groupe de $\mathcal M_{0,3}^{!\sim}$
formé des $u$ qui satisfont aux conditions équivalentes (93) (94) (98)
\[
 (100)\qquad
 \mathcal M_{0,3}^{!\sim}\left(\tfrac12\right)
 =
 \{u\in\mathcal M_{0,3}^{!\sim}\mid [u,\sigma_\infty]=1\}
 \overset{\sim}{\longrightarrow}\mathcal N_{0,3}^{!\sim}
\]
on trouve que l'hom composé

\scanpage{208}
\mspage{403}

(101)
\diagram{
\begin{tikzcd}[ampersand replacement=\&,column sep=large]
\mathcal N_{0,3}^{!\sim}
 \arrow[r,"\sim"]
 \arrow[rr,bend right=22,"k_{1/2}"']
\&
\mathcal M_{0,3}^{!}(\tfrac12)^\sim
 \arrow[r]
\&
\mathcal M_{0,3}^{!\sim}
 \arrow[r,"{u\longmapsto\underline{u}=(u_{\alpha,\beta,f})}"']
\&
\mathrm{Gl}(2,\widehat{\mathbb Z})'
\\
\unsure{\mathcal N_{0,3}^{!\sim}(\tfrac12)}
\end{tikzcd}
}

\margin{\uline{NB} $\mathcal M_{0,3}^{!}(\tfrac12)^\sim$
est stable par $x\longmapsto\sigma_\infty x\sigma_\infty^{-1}$
et par $x\longmapsto\tau_\infty x\tau_\infty^{-1}$
(puisque $\sigma_\infty$ et $\tau_\infty$
\unsure{commutent} à $\sigma_\infty$).}

factorise par le dit normalisateur --- si on préfère le voir en termes du
groupe $\mathrm{Gl}(2,\widehat{\mathbb Z})\simeq M_2(\widehat{\mathbb Z})^*$,
plutôt que de l'\uline{algèbre} des matrices, elle se factorise par
le \uline{normalisateur} du \uline{tore maximal}
\deleted{(\unsure{formé par})}
\[
 (102)\qquad
 \unsure{T_\alpha|T_\sigma}=
 \{t\sigma_\infty+u\quad
 (t,u\in\unsure{\widehat{\mathbb Z}|\mathbb Z},
 \underbrace{t^2+u^2}_{=\det(t\sigma_\infty+u)}
 \in\widehat{\mathbb Z}^{*}).
\]

Je suspecte qu'il n'y a \uline{que le} scindage $k_{1/2}$ de (101)
de l'extension $\mathcal M_{0,3}^{!\sim}$
(de $\mathcal N_{0,3}^{!\sim}$ par $\widehat{\Pi}_{0,3}$,
\deleted{\unsure{Je crois}} ou de $\mathcal N_{0,3}^{!\sim}$
par $\widehat{\mathcal G}_{0,3}^{+}$)
qui se factorise par le normalisateur.

\margin{Sûrement canulé :
il y a aussi \unsure{toutes} les courbes \uuline{\unsure{isogènes}}
à la courbe elliptique d'invariant modulaire $\unsure{j|J}=0$
(\unsure{correspondant} \deleted{dans}
l'orbite $\{\tfrac12,2,-1\}$ dans $U_{0,3}$ sous $\unsure{\mathfrak S_3}$).}

\uline{2° Cas}\quad C'est le cas
\[
 (103)\qquad
 g=1\quad\text{i.e. }\unsure{\alpha_0}=1
 \quad\text{i.e. }\alpha=f
 \quad\text{i.e. }\xi=f\rho(f)^{-1},\ \text{i.e.}
\]
\[
 (104)\qquad [u,\rho]=1.
\]

Pour un $u$ quelconque, il y a un
$\gamma\inserted{\in\widetilde{\widehat{\Pi}}_{0,3}
=\widehat{\Pi}_{0,3}\cdot\{1,\tau_\infty\}}$
\deleted{\unsure{unique}} tel que
$u'=\mathrm{int}(\gamma)\circ u=\gamma u$
\unsure{commute|commutant} ainsi à $\rho$,
savoir $\gamma=\unsure{\alpha_0^{-1}}$,
\deleted{\unsure{On a donc}}\inserted{Mais} on fera Attention que
$\gamma$ n'est pas toujours dans $\widehat{\Pi}_{0,3}$, en fait
\[
 (105)\qquad
 \gamma=\unsure{\alpha_0^{-1}}\in\widehat{\Pi}_{0,3}
 \Longleftrightarrow
 \varepsilon_3(\unsure{u})=1
 \quad\text{i.e. }\mu(u)=1\ (3)
\]

\scanpage{209}

Si on désigne par
$\unsure{\mathcal M_{0,3}^{!\sim}}(\unsure{-\overline{j}})$
le sous-groupe de $\mathcal M_{0,3}^{!\sim}$
formé des $u$ satisfaisant (103), (104), son image dans
$\unsure{\mathcal N_{0,3}^{!\sim}}$ est donc
le sous-groupe d'indice $2$ noyau de
\[
 u\longmapsto
 \underbrace{\varepsilon_3(\mu(u))}_{\unsure{\sim}}
 =\varepsilon_3(u)
\]
i.e.
\[
 \unsure{(106)}\qquad
 \mathcal N_{0,3}^{!\sim}(-\overline{j})
 =\{u\in\mathcal N_{0,3}^{!\sim}\mid
 \varepsilon_3(u)=1\quad
 \text{i.e. }\mu(u)\equiv1\ (3)\}
\]

\margin{\uline{NB} $\mathcal M_{0,3}^{!\sim}(-\overline{j})$
est stable par $x\longmapsto\rho x\rho^{-1}$.}

et l'on a
\[
 \unsure{(107)}\qquad
 \unsure{\mathcal M_{0,3}^{!\sim}}[-\overline{j}]
 \overset{\sim}{\longrightarrow}
 \mathcal N_{0,3}^{!\sim}[-\overline{j}].
\]

En l'inversant on trouve donc un
\deleted{hom composé}\inserted{scindage}

(108)
\diagram{
\begin{tikzcd}[ampersand replacement=\&,column sep=large]
\mathcal N_{0,3}^{!\sim}(\unsure{-\overline{j}})
 \arrow[r,"\sim"]
 \arrow[rr,bend right=20,"{k(-\overline{j})}"']
\&
\mathcal M_{0,3}^{!\sim}(\unsure{-\overline{j}})
 \arrow[r,hook]
\&
\mathcal M_{0,3}^{!\sim}
\end{tikzcd}
}
,
dont le composé avec $\mathcal M_{0,3}^{!\sim}$ est un hom qui
\deleted{\unsure{factorise}}\inserted{se factorise}
par le tore maximal $T_\rho$ de $\mathrm{Gl}(2,\widehat{\mathbb Z})'$,
formé des
\deleted{$\unsure{c\rho+d}$}
\[
 \unsure{(109)}\qquad
 \{c\rho+d\mid c,d\in\widehat{\mathbb Z},\
 c^2+cd+d^2\in\widehat{\mathbb Z}^{*}).
\]

Je suspecte que c'est le seul scindage de
$\mathcal M_{0,3}^{!\sim}$
(ou \unsure{$\overline{\mathrm m}$} de $\mathcal M_{0,3}^{!}$)
qui ait cette propriété. \unsure{Celle-ci} résulte de la forme équivalente
des conditions (103) (104), fournie par (79) :

\margin{\unsure{canulé}, pour \unsure{mêmes} raisons que tout à l'heure\ldots}

\scanpage{210}
\mspage{404}

\[
 (110)\qquad
 \underline{u}=u_{\alpha,\beta,f}=\mu(c\rho+d)
\]
(NB $\delta=\varepsilon_3(\unsure{\mu|u})$ est égal à $1$ dans ce cas).

Les $u$ en question correspondent aux couples $(\eta,f)$ satisfaisant
\[
 (111)\qquad \eta\sigma_\infty(\eta)=1
\]
\[
 \unsure{(112)}\qquad
 (f\rho(f)^{-1})=\eta\ell_1^{\unsure{\nu}}
\]

\margin{qu'on \unsure{complètera|complète} par $\xi=f\rho(f)^{-1}$.}

\margin{On peut aussi écrire $\beta$ au lieu de \unsure{$f$},
$f(=\beta)$ \unsure{déterminé} mod
$f\longmapsto\ell_0^{\unsure{n}}f$,
\unsure{id.} pour $\alpha,\beta$.}

ou encore aux couples $(\alpha,\beta)$ satisfaisant la relation des lacets
$(97)=(37)$, qu'on complète alors par $f=\alpha$.

\uline{3° Cas}\quad C'est celui où
\[
 (113)\qquad
 f\equiv1\ (L_0)
 \quad\text{i.e. }\alpha\equiv\alpha_0\ (L_0\text{ à g.})
 \quad\text{i.e. }\beta\equiv\beta_0\ (L_0\text{ à g.})
\]
ou encore $u(\varepsilon_0)=\varepsilon_0^\mu$
$(\mu\in\widehat{\mathbb Z}^{*})$ i.e.
\[
 (114)\qquad
 [u,\varepsilon_0]=\ell_0^\nu
 \qquad
 (\nu\in\widehat{\mathbb Z}\ -\ \nu=\tfrac{\mu-1}{2}),
\]
\inserted{i.e. $u$ normalise
$L_0=\{\ell_0^n\mid n\in\widehat{\mathbb Z}\}$ --- ou encore
(\unsure{agrandissant} $\widehat{\Pi}_{0,3}$)
$u$ normalise $L'_0=\{\varepsilon_0^n\mid n\in\widehat{\mathbb Z}\}$.}

\margin{\uline{NB} $L_0=L'_0\cap\widehat{\Pi}_{0,3}$.}

Dans ce cas, on \unsure{lève|peut lever} l'indétermination
$(\alpha,\beta,f)\longmapsto
(\ell_0^n\alpha,\ell_0^n\beta,\ell_0^n f)$
en exigeant
\[
 (115)\qquad f=1,
\]
ce qui détermine $\alpha,\beta$ de façon unique.

\deleted{L'ensemble}\inserted{groupe}
$\mathcal M_{0,3}^{!\sim}[0]$ des $u\in\mathcal M_{0,3}^{!\sim}$
qui normalisent $L_0$ est en corr. biunivoque avec l'ens des couples
$(\alpha,\beta)$
$(\alpha\in\widetilde{\widehat{\Pi}}_{0,3}
=\widehat{\Pi}_{0,3}\cdot\{1,\tau_\infty\},
\ \beta\in\widehat{\Pi}_{0,3})$ satisfaisant

\margin{On ne \unsure{remonte|relève} Donc \illegible{2}
par $\mathcal M_{0,3}^{!\sim}(0)$,
\illegible{2} $\mathcal M_{0,3}$ et $\mathcal N_{0,3}$
car il \unsure{ne couvre} \illegible{2}
section de l'ext. \illegible{2} par $\widehat{\Pi}_{0,3}$
--- on \illegible{2} (Prop. \illegible{1})
section $\mathcal M_{0,3}^{!\sim}(0)$ plus
\unsure{p. 412$^1$}).}

\scanpage{211}

la sempiternelle relation des lacets
\[
 \unsure{(115)}\qquad
 (\alpha\rho(\alpha)^{-1})
 =(\beta\sigma_\infty(\beta)^{-1})\ell_1^\nu
 \qquad(\text{pour }\nu\in\widehat{\mathbb Z}\text{ conv.})
\]
\inserted{(par $(\alpha,\beta)\longmapsto
\unsure{g|\xi}=\alpha\rho(\alpha)^{-1};
\unsure{h|\eta}=\beta\sigma_\infty(\beta)^{-1}$)}
ou encore aux couples $(g,h)$
\inserted{(\unsure{ou} $(\xi,\eta)$)}
d'éléments de $\widehat{\Pi}_{0,3}$ satisfaisant
\[
 \begin{gathered}
 (116)\quad (=(44)+(46))\\
 g\rho(g)\rho^2(g)=1,\qquad
 \unsure{h}\sigma_\infty(\unsure{h})=1,\qquad
 g=\unsure{h}\ell_1^\nu\quad
 (\nu\in\widehat{\mathbb Z}\text{ conv.}),
 \end{gathered}
\]

\margin{Dans tout ceci, je laisse tomber la "condition
liminaire" sur les relations génératrices\ldots}

au couple \deleted{$(\alpha,\beta)$ \unsure{comme}}
$(g,h)$ correspond l'automorphisme $u$ défini par
\[
 (117)=(23)\qquad
 u(\sigma_\infty)=\unsure{h}\sigma_\infty,\qquad
 u(\rho)=g\rho
\]
ou encore par
\[
 (118)\qquad
 \unsure{u}(\sigma_\infty)=\mathrm{int}(\alpha)\sigma_\infty,\qquad
 u(\rho)=\mathrm{int}(\beta)\rho.
\]

Dans le groupe $\mathcal M_{0,3}^{!\sim}[0]$, il y a le sous groupe
\[
 (119)\qquad
 \mathcal M_{0,3}^{!\sim}[0]\cap\widehat{\Pi}_{0,3}\simeq L_0
\]
l'inclusion
\[
 (120)\qquad
 \widehat{\mathbb Z}\overset{\ell_0}{\underset{\sim}{\longrightarrow}}
 L_0\lhook\joinrel\longrightarrow\mathcal M_{0,3}^{!\sim}
\]
étant décrite par
\[
 (121)\qquad
 \lambda=\ell_0^n\longmapsto
 (\alpha=\unsure{\lambda|\ell_0^n},
 \beta=\unsure{\lambda|\ell_0^n})
\]
\deleted{\unsure{i.e.}\ \illegible{2}\ $(\alpha=\lambda,\beta=$}

ou encore par
\[
 (122)\qquad
 \lambda=\ell_0^{\unsure{n}}\longmapsto
 (f=\ell_0^n\ell_1^{-n},g=\ell_0^n\ell_1^{-n})
\]
\inserted{i.e. $f=g=\lambda\sigma_\infty(\unsure{\lambda^{-1}})
=\lambda\rho(\unsure{\lambda})^{-1}$,}
la multiplication \deleted{par} à g. par \unsure{le}
$\lambda=\ell_0^n$ dans $\mathcal M_{0,3}^{!\sim}[0]$
étant \unsure{explicité|explicitée} par
\[
 (123)\qquad
 (\alpha,\beta)\longmapsto(\ell_0^n\alpha,\ell_0^n\beta)
\]
\[
 (124)\qquad
 (\xi,\eta)\longmapsto
 \ell_0^n\xi\ell_1^{-n},\ \ell_0^n\eta\ell_1^{-n}).
\]

\scanpage{212}
\mspage{405}

Comme on a un iso
\[
\mathcal M_{0,3}^{!\sim}[0]/L_0
\xrightarrow{\sim}\mathcal N_{0,3}^{!\sim},
\]
on trouve le

\uline{Corollaire} Le groupe $\mathcal N_{0,3}^{!\sim}$ est en bijection
canonique avec l'ens des couples
$(\alpha,\beta)\in\Pi_{0,3}^{\sim\wedge}\times\Pi_{0,3}^{\wedge}$
satisfaisant la relation des lacets (97)
\[
\alpha\rho(\alpha)^{-1}
=\bigl(\beta\sigma_\infty(\beta)^{-1}\bigr)\ell_1^\nu
\qquad(\text{pour }\nu\in\widehat{\mathbb Z}\text{ conv.}),
\]
\uline{modulo} les transformations
$(\alpha,\beta)\mapsto(\ell_0^{\unsure{n|m}}\alpha,\ell_0^{\unsure{n|m}}\beta)$.

La loi de groupe \inserted{(de $\mathcal M_{0,3}^{!\sim}(0)$ ou de son quotient
$\mathcal N_{0,3}^{!\sim}$)} s'explicite en termes des $(g,h)$, ou des
$(\alpha,\beta)$, par les ``formules des cocycles'' (52) : si $u',u$
\inserted{et $u''=u'\cdot u$} sont donnés par
$\unsure{(g',h')},(g,h),(g'',h'')$, ou par
$(\alpha',\beta'),(\alpha,\beta)$ et $(\alpha'',\beta'')$, on aura
\[
\tag{125}
\left\{
\begin{aligned}
g''&=u'(g)g'&h''&=u'(h)h'\\
\alpha''&=u'(\alpha)\alpha'&\beta''&=u'(\beta)\beta'.
\end{aligned}
\right.
\]

\uline{Remarque sur les notations} $\mathcal M_{0,3}^{!\sim}(q)$,
pour $q=\frac12,-\overline j,0$.
En fait pour tout point \uline{algébrique}
\[
\tag{126}
q\in\mathcal U_{0,3}(\mathbb C)
=\mathbb C\setminus\{0,1\}
=\mathbb P^1_{\unsure{\mathbb C|\mathbb Q}}(\mathbb C)
\setminus\{0,1,\infty\}
\]
on s'attend à \deleted{\unsure{trouver}} un scind. partiel
\[
\tag{127}
k_q:\unsure{\mathcal N|\mathcal M}_{0,3}^{!\sim}(q)\longrightarrow\mathcal M_{0,3}^{!\sim},
\qquad\operatorname{Im}k_q\overset{\text{déf}}=\mathcal M_{0,3}^{!\sim}(q)
\]
de l'extension $\mathcal M_{0,3}^{!\sim}$ de
$\mathcal N_{0,3}^{!\sim}$ par $\Pi_{0,3}^{\wedge}$, ou

\margin{En fait, \unsure{de façon} \illegible{1} plus
\unsure{précise|précises}, $k_q$ ne sera \unsure{déterminé}
qu'à \unsure{autom.} int. près...}

\scanpage{213}

du moins \inserted{on aura} un scindage partiel au dessus
de $\Gamma_{\mathbb Q}(q)$, qu'on s'efforcera de déterminer ;
$\mathcal N_{0,3}^{!\sim}(q)$ désigne un sous-groupe
hypothétique de $\mathcal N_{0,3}^{!\sim}$
(qu'il faudrait \unsure{arriver} à déterminer, \unsure{si tant est}
\illegible{1} pas trop \illegible{1}...), induisant sur $\Gamma_{\mathbb Q}$
le sous-groupe \inserted{ouvert} $\Gamma_{\mathbb Q}(q)$ de
$\Gamma_{\mathbb Q}$ stabilisateur de $q$.

\margin{La définition s'impose p.ex. lorsque $q$ est
\unsure{élément} d'une extension cyclotomique de $\mathbb Q$, de telle
façon que $\mathbb Q(q)$ correspond à un groupe quotient fini
discret de $\widehat{\mathbb Z}^{*}$ et $\Gamma_{\mathbb Q}(q)$ est
le noyau de
\diagram{\begin{tikzcd}[ampersand replacement=\&]
\Gamma_{\mathbb Q}\arrow[r,"\chi"]\&\widehat{\mathbb Z}^{*}\arrow[r]\&G
\end{tikzcd}}\ldots}

L'analyse des cas $1^\circ$ et $2^\circ$ ci-dessus revenait en fait à
la détermination de tels scindages, pour $q=\frac12$ et pour
$q=-\overline j$ --- dans le premier cas $\Gamma_{\mathbb Q}(q)$ est
$\Gamma_{\mathbb Q}$ tout entier, dans le deuxième c'est le sous-groupe
d'indice $2$ noyau de $\varepsilon_2\circ\chi:\Gamma_{\mathbb Q}
\longrightarrow\{\pm1\}$, et idem pour
$\unsure{\mathcal N|\mathcal M}_{0,3}^{!\sim}(q)$.
En transformant par $\rho$ \inserted{et $\rho^2$}, on
\deleted{\unsure{conduirait}}\inserted{déduirait} \unsure{des scindages|du scindage} $k_{1/2}$
les scindages $k_2$ et $k_{-1}$ --- en transformant par un quelconque
des $\sigma_{\unsure{i|\rho}}$,
\deleted{ou par un des $\unsure{\tau_i|\sigma_i}$}, on déduirait
du scindage $k_{-\overline j}$ un scindage du type $k_{\unsure{-j|-1}}$.

Quant à $\mathcal M_{0,3}^{!\sim}[0]$, il correspond
\inserted{en principe} au point

\scanpage{214}
\mspage{406}

$0$ de $\widehat{\mathcal U}_{0,3}=\unsure{\mathbb P^1_{\mathbb Q}}$,
il \deleted{\unsure{y correspondent par}}\inserted{deux autres}
sous-groupes analogues $\mathcal M_{0,3}^{!\sim}[q]$, avec
$q=1$, $q=\infty$, correspondant aux deux autres points ``à l'infini''
de $\mathcal U_{0,3}$, et qui se déduisent du précédent en appliquant
$\rho$, $\rho^2$.
Ici, on a des groupes qui \uline{ne} correspondent \uline{pas}
\inserted{\unsure{strictement}} à des sections \inserted{même} partielles
de $\mathcal M_{0,3}^{!\sim}$ sur $\mathcal N_{0,3}^{!\sim}$,
puisque l'intersection avec $\Pi_{0,3}^{\wedge}$ n'est pas réduite à $1$,
mais c'est l'un des trois groupes-lacets type $L_0,L_1,L_\infty$.

Mais on sait que \inserted{le choix d'}un \unsure{uniformisant|uniformisante}
de l'anneau local de $\widehat{\mathcal U}_{0,3}$ en $0$ (disons)
détermine un scindage de
$\mathcal M_{0,3}^{!\unsure{\sim}}\unsure{(0)|[0]}$, considéré comme
ext. de $\mathcal N_{0,3}^{!}$ par $L_0$ --- et il faudra
dans la suite essayer de saisir \deleted{\unsure{ce}}\inserted{un tel}
scindage \inserted{(associé à une \unsure{uniformisante})}
dans l'optique de notre calcul, comme une façon \unsure{privilégiée}
\deleted{de ``\unsure{normaliser}''}, pour \unsure{tte}
\deleted{$\alpha,\beta$} \unsure{donnée} \inserted{mod $L_0$} de $(\alpha,\beta)$
satisfaisant une relation des lacets, de ``normaliser'' le choix des
$(\alpha,\beta)$ mod $L_0$.
J'y reviendrai par la suite, sûrement.

\scanpage{215}
\mspage{407}

\uline{V} \uline{Récapitulation} \uline{en termes de}
$\mathcal M_{0,3}^{!\sim}[0]$

Il apparaît finalement que la façon la plus simple d'étudier les autom.
à lacets de $\Pi_{0,3}^{\wedge}$ (commutant ext.\textsuperscript{t}
à $\mathfrak S_3$, i.e. les éléments de \deleted{\illegible{1}}
$\mathcal N_{0,3}^{!\unsure{\sim}}$) est de commencer par l'étude du
sous-groupe $\mathcal M_{0,3}^{!\unsure{\sim}}[0]$, normalisateur
de $L_0$ \inserted{dans $\mathcal M_{0,3}^{!\sim}$}.
Posant pour $u\in\mathcal M_{0,3}^{!\unsure{\sim}}[0]$
\[
\tag{128}
u(\sigma_\infty)=\eta\sigma_\infty,
\qquad u(\rho)=\xi\rho
\qquad(\eta,\xi\in\Pi_{0,3}^{\wedge})
\]
i.e.
\[
\tag{129}
[u,\sigma_\infty]=\eta\qquad [u,\rho]=\xi
\]
qui expriment la commutation ext. de $u$ à $\mathfrak S_3$
(engendré en effet par $\sigma_\infty$ et $\rho$), la condition de
normalisation de $L_0$ s'exprime par une relation de forme
\[
\tag{130}
u(\varepsilon_0)=\varepsilon_0^\mu
\quad\text{ou}\quad \unsure{u}(\ell_0)=\ell_0^\mu
\qquad(\mu\in\widehat{\mathbb Z}^{*}).
\]

\margin{(130) équivaut aussi, moyennant $[u,\rho]=\xi$, à
\[
\tag{130bis}
\left\{
\begin{aligned}
u(\varepsilon_1)&=\operatorname{int}(\xi)(\varepsilon_1^\mu)\\
&\text{et à}\\
u(\varepsilon_\infty)&=
\operatorname{int}\bigl(\rho^{-1}(\xi^{-1})\,\illegible{1}\bigr)(\varepsilon_\infty^\mu)
\end{aligned}
\right.
\]}

\inserted{Or on a, comme $\varepsilon_0=\sigma_\infty\rho$}
\deleted{qui, compte tenu de $\varepsilon_0=\sigma_\infty\rho$,
prend (par un calcul immédiat) la forme}
\deleted{\unsure{s'écrit}}
\deleted{(131)} \deleted{$\unsure{\eta\sigma_\infty}$}
\[
u(\varepsilon_0)=u(\sigma_\infty)u(\rho)
=\eta\sigma_\infty\xi\rho
=\eta\sigma_\infty(\xi)\underbrace{\sigma_\infty\rho}_{\varepsilon_0}
=\eta\sigma_\infty(\xi)\varepsilon_0
\]
donc (130) s'écrit
\[
\eta\sigma_\infty(\xi)=\varepsilon_0^{\mu-1}
\qquad\text{soit, en posant }\mu-1=2\nu
\quad(\nu\in\widehat{\mathbb Z})
\]
\[
\eta\sigma_\infty(\xi)=\ell_0^\nu
\qquad\text{soit, appliquant }\sigma_\infty
\]
\[
\tag{131}
\sigma_\infty(\eta)\xi=\ell_1^\nu.
\]
Or les relations
\[
\tag{132}
(\eta\sigma_\infty)^2=1,\qquad(\xi\rho)^3=1
\]
s'écrivent aussi
\[
\tag{133}
\boxed{\eta\sigma_\infty(\eta)=1,\qquad
\xi\rho(\xi)\rho^2(\xi)=1},
\]
d'où $\sigma_\infty(\eta)=\eta^{-1}$, de sorte que (131) s'écrit
aussi de façon plus \uline{commode}
\[
\tag{134}
\boxed{\xi=\eta\ell_1^\nu}
\qquad(\text{relation des lacets})
\]

\scanpage{216}

On trouve ainsi le

\uline{Théorème} Les éléments $u$ de
$\unsure{\mathcal M|\mathcal N}_{0,3}^{!\sim}[0]$ --- i.e. les
autom. ext. de $\Pi_{0,3}^{\wedge}$ qui commutent ext.\textsuperscript{t}
à $\mathfrak S_3$ et qui normalisent $L_0$ --- sont en corr.
biunivoque \inserted{(via les formules (128) ou 129)}
avec les couples $(\xi,\eta)$ $(\xi,\eta\in\Pi_{0,3}^{\wedge})$,
satisfaisant les relations (133)(134), et tels que
\inserted{de plus (``condition liminaire'')}
les relations (132) (équivalentes : (133)) soient
\deleted{\unsure{primitives}}\inserted{génératrices}
pour $\widehat{\mathcal G}_{0,3}^{+}$.
(\deleted{\illegible{1}} Dans (134) $\nu$ est un élément convenable de
$\widehat{\mathbb Z}$, évidemment uniquement déterminé, et assujetti
à la condition supplémentaire $2\nu+1\in\widehat{\mathbb Z}^{*}$)
(qui est impliquée par la condition liminaire, et peut-être l'implique...)

\margin{cette condition $2\nu+1\in\widehat{\mathbb Z}^{*}$ est en fait
conséquence des équations (133), (134), cf plus bas th. p. 409!}

\unsure{Et} l'inclusion évidente
\[
\tag{135}
\begin{aligned}
L_0&\hookrightarrow\mathcal M_{0,3}^{!\unsure{\sim}}\unsure{(0)|(\omega)}\\
\ell_0^n=\lambda&\longmapsto\operatorname{int}(\lambda)
\end{aligned}
\]
est donnée par
\[
\tag{136}
\ell_0^n=\lambda\longmapsto
\bigl(\xi=\lambda\rho(\lambda)^{-1}=\ell_0^n\ell_1^{-n},
\quad\eta=\lambda\sigma_\infty(\lambda)^{-1}=\ell_0^n\ell_1^{-n}\bigr)
\]
et plus généralement l'action à gauche de $L_0$ sur
$\mathcal M_{0,3}^{!\unsure{\sim}}\unsure{(0)|(\omega)}$ est donnée,
pour $\lambda=\ell_0^n\in L_0$, par
\[
\tag{137}
(\xi,\eta)\longmapsto
\bigl(\xi'=\ell_0^n\xi\ell_1^{-n},\quad
\eta'=\ell_0^n\eta\ell_1^{-n}\bigr)
\]
donc
\[
\mathcal M_{0,3}^{!\unsure{\sim}}[0]/L_0
\simeq\mathcal N_{0,3}^{!\unsure{\sim}}
\]
est isomorphe au quotient de l'ensemble des couples $(\xi,\eta)$
satisfaisant (133)(134) [et la condition liminaire], par les opérations
(137) --- compte tenu d'une \deleted{\unsure{opération}}
structure de groupe naturelle : mettre sur cet ens des couples
$(\xi,\eta)$, donnée par

\scanpage{217}
\mspage{408}

\[
\tag{138}
\underbrace{(\xi',\eta')}_{\deleted{\unsure{u'}}}(\xi,\eta)=(\xi'',\eta'')
\]
avec
\[
\tag{139}
\xi''=u'(\xi)\xi',\qquad \eta''=u'(\eta)\eta'
\]
$u'$ étant l'automorphisme défini par $(\xi',\eta')$.

\uline{Corollaire} Tout $u\in\unsure{\mathcal M_{0,3}^{!\sim}}$ s'écrit \unsure{sous|dans} la forme
\[
\tag{140}
u=f^{-1}u_{\xi,\eta}=\operatorname{int}(f^{-1})\circ u_{\xi,\eta},
\]
avec $f\in\widehat{\Pi}_{0,3}$ et $(\xi,\eta)$ satisfaisant aux conditions du théorème précédent, de sorte qu'on aura
\[
\tag{141}
u(\unsure{\sigma_\infty})=h\sigma_\infty,\qquad u(\rho)=g\rho,\quad\text{avec}
\]
\[
\tag{143}
h=f^{-1}\eta\sigma_\infty(f),\qquad g=f^{-1}\xi\rho(f).
\]
Pour qu'on ait
\[
f^{-1}u_{\xi,\eta}=f'^{-1}u_{\xi',\eta'}
\]
il \unsure{f. et s.} qu'il existe $n\in\mathbb Z$ tel que
\[
\tag{144}
\xi'=\ell_0^n\xi\ell_1^{-n},\qquad
\eta'=\ell_0^n\eta\ell_1^{-n},\qquad f'=\ell_0^n f.
\]

\uline{NB} \unsure{Grâce} à la formule des lacets, une fois fixé le multiplicateur $\mu$ i.e. $\nu$ (et en fait $\mu$ est déjà déterminé de façon unique par la seule connaissance de $\xi$, comme le \unsure{montrent} les calculs adéliques dans $\mathrm{Gl}(2,\widehat{\mathbb Z})'\ldots$), $\xi$ et $\eta$ se déterminent mutuellement \inserted{par (134)}, donc on peut éliminer l'un ou l'autre du système d'équations (133), (134), qui devient soit
\[
\tag{145}
\begin{gathered}
\xi\rho(\xi)\rho^2(\xi)=1,\qquad
(\deleted{\unsure{\eta}}\inserted{\xi}\ell_1^{-\nu})\sigma_\infty(\xi\ell_1^{-\nu})=1
\quad\text{i.e.}\\
\sigma_\infty(\xi)=\ell_1^\nu\xi^{-1}\ell_0^\nu
\end{gathered}
\]
soit
\[
\tag{146}
\eta\sigma_\infty(\eta)=1,\qquad
(\eta\ell_1^\nu)\rho(\eta\ell_1^\nu)\rho^2(\eta\ell_1^{\unsure{\nu|-\nu}})=1.
\]

\scanpage{218}

Le calcul adélique. On désigne encore par $\xi,\eta$ les images de $\xi,\eta$ dans $\mathrm{Sl}(2,\widehat{\mathbb Z})'$, qui remontent en des éléments, notés encore $\xi,\eta$, de $\mathrm{Sl}(2,\widehat{\mathbb Z})$.

Il y a une ambiguïté de signe, qu'on lève pour $\xi$ en imposant
\[
\tag{147}
\xi\rho(\xi)\rho^2(\xi)=\uuline{+}1
\qquad(\text{dans }\mathrm{Sl}(2,\widehat{\mathbb Z}))
\]
\inserted{pour $\eta$} il y \uline{a pas de façon simple} de lever l'\unsure{ambiguïté} (si on veut admettre \inserted{provisoirement} la possibilité d'écrire sous la forme $\beta\sigma_\infty(\beta)^{-1}$) \uline{\unsure{dont|donc}} on aura a priori seulement
\[
\tag{148}
\eta\sigma_\infty(\eta)=\deleted{\unsure{\varepsilon}}\pm1,
\]
sans possibilité de modifier le signe au $2^\circ$ cas par changement $\eta\mapsto-\eta$ ! Posons

\margin{\uline{mais si}, on peut faire le choix de $\eta$ en exigeant $\xi=+\eta\ell_1^\nu$ (mais cette convention \unsure{n'est} pas compatible avec une convention \unsure{faite} par ailleurs $\eta=\beta\sigma_\infty(\beta)^{-1}$).}
\margin{$\in\widehat{\Pi}_{0,3}$}

\[
\tag{148}
\xi=\begin{pmatrix}A&B\\C&D\end{pmatrix},\qquad
\eta=\begin{pmatrix}R&S\\T&U\end{pmatrix}.
\]
On aura
\[
\tag{149}
AD-BC=1\qquad RU-ST=1
\]
et la relation (147) équivaut alors à :
\[
\tag{150}
C+\unsure{D}-B=1,
\]
\deleted{\unsure{tandis}} (du moins dans chaque composante $\ell$-adique on devra avoir soit cette relation, soit
\[
\xi_\ell=\begin{pmatrix}A_\ell&B_\ell\\C_\ell&D_\ell\end{pmatrix}
=-\rho^{-1}=\begin{pmatrix}1&-1\\-1&0\end{pmatrix},
\]
\deleted{\unsure{cf.}} prop. de la \unsure{p.g} 394 -- mais cette solution est incompatible avec la condition $\boxed{\xi\equiv1\ (2)}$, d'où $A,D\equiv1\ (2)$ et $B,C\equiv0\ (2)\ldots$ donc (150) est vraie dans $\mathbb Z_\ell$ pour $\forall\ell$, donc dans $\widehat{\mathbb Z}\ldots$).

\scanpage{219}
\mspage{409}

La relation (148) s'écrit aussi
\[
\eta=\pm\sigma_\infty(\eta^{-1})
\]
or $\sigma_\infty(\eta^{-1})={}^t\eta$, donc (148) s'écrit
\[
\eta=\pm{}^t\eta,
\]
mais si on avait le signe $-$, on aurait $U=-U$ donc $2U=0$ donc $U=0$, ce qui est incompatible avec
\[
\eta\equiv\unsure{\pm1|1}\quad(2)
\]
qui implique $R,U\equiv1\ (2)$, $S,T\equiv0\ (2)$. On aura donc $\eta={}^t\eta$ i.e.
\[
\tag{151}
\eta\sigma_\infty(\eta)=+1,\qquad\text{i.e.}
\]
\[
\tag{152}
S=T.
\]

La relation des lacets (134) s'écrit
\[
\tag{153}
\begin{aligned}
\begin{pmatrix}A&B\\C&D\end{pmatrix}
&=\deleted{\begin{pmatrix}1&0\\2\nu&1\end{pmatrix}}
\inserted{\varepsilon}\begin{pmatrix}R&S\\T&U\end{pmatrix}
\begin{pmatrix}1&0\\2\nu&1\end{pmatrix}\\
&=\varepsilon\begin{pmatrix}R+2\nu S&S\\T+2\nu U&U\end{pmatrix}
\qquad\text{avec }\varepsilon\in\{\pm1\}
\end{aligned}
\]

\margin{\unsure{Obs} ici $\varepsilon=1$, en choisissant le signe de $\eta$ convenablement (153) (la convention de signe ne coïncide pas avec une convention antérieure axée sur les $\alpha,\beta$, mais dans l'optique présente elle paraît plus \unsure{sympa}).}

et elle équivaut aux relations
\[
\tag{154}
B=\varepsilon S,\qquad D=\varepsilon U\qquad\text{et}
\]
\deleted{\illegible{1}}
\[
\tag{155}
\mu D=1\qquad\text{où }\mu=2\nu+1
\]
\unsure{équivalente|équivalant} aussi ; moyennant (154), à
\[
\tag{155bis}
\mu U=\varepsilon,
\]
et cette dernière ne fait qu'exprimer \inserted{(compte tenu de (152))} la relation
\[
\varepsilon\overbrace{(T+2\nu U)}^{A'}
+\overbrace{\varepsilon U}^{B'}-\overbrace{\varepsilon S}^{C'}=1
\]
\unsure{paraphrasant}
\[
A+B-C=1,
\]
elle

\scanpage{220}

elle ne fait qu'exprimer (compte tenu de (154) et de \unsure{(150)} la relation
\[
C=\overbrace{\varepsilon(T+2\nu U)}^{C'}.
\]
Ainsi les relations (154) (155) expriment le fait que les deux membres de (153) \deleted{\unsure{tous} \unsure{leurs termes}} \inserted{trois coefficients} égaux \deleted{\illegible{1}} \inserted{\unsure{deux à deux}}, \inserted{$\unsure{B=B'},\ C=C',\ D=D'$}, et l'égalité pour le dernier
\[
A=\overbrace{\varepsilon(R+2\nu S)}^{A'}
\]
s'en déduit alors grâce aux relations \deleted{(149)}, \deleted{qui permettent d'exprimer}
\[
AD-BC=1,\qquad A'D'-B'C'=1,
\]
qui donnent
\[
A=\frac{1+BC}{D},\qquad A'=\frac{1+B'C'}{D'}
\]
(NB on sait que $D=D'$ est inv., par (155)) donc $A=A'$.

On trouve ainsi le

\uline{Théorème} Pour une solution $(\xi,\eta)$ dans $\widehat{\Pi}_{0,3}\times\widehat{\Pi}_{0,3}$ du système d'équations (133) (134), désignant encore par $\xi,\eta$ des éléments de $\mathrm{Sl}(2,\widehat{\mathbb Z})$ qui leur correspondent, l'ambiguïté de signe pour $\xi$ étant levée par la relation (147), on aura les relations (154) et (155), qui impliquent les conséquences suivantes :

a) $\mu\inserted{=2\nu+1}$ est inversible, et il est entièrement \unsure{déterminé} par $\xi$ -- il est déterminé par $\eta$ seulement au signe près.

\scanpage{221}
\mspage{410}

\uline{Remarque}

Compte tenu des relations (149) (150)
\[
AD-BC=1,\qquad \unsure{A|C}+D-B=1,
\]
et du fait que $D$ est inv., les matrices $\xi$ sont entièrement déterminées par les coeff $B,D$ --- les couples $(\xi,\eta)$ de matrices $\in\mathrm{Sl}(2,\widehat{\mathbb Z})$, satisfaisant
\[
\tag{156}
\left\{\begin{gathered}
\xi\rho(\xi)\rho^2(\xi)=1,\quad
\eta\sigma_\infty(\eta)=1,\quad
\xi=\eta\ell_1^\nu
\qquad(\nu\in\widehat{\mathbb Z}\ \text{\unsure{convenant}})\\
\xi\equiv1\quad(2),\qquad\eta\equiv1\quad(2)
\end{gathered}\right.
\]

\margin{Bien sûr, l'introduction des congruences ici est \unsure{tout à fait} accessoire --- l'énoncé algébrique \unsure{est indépendant} de telles conditions...}

sont en correspondance biunivoque avec les couples $(B,D)$, $B\in\widehat{\mathbb Z}$, $D\in\widehat{\mathbb Z}^{*}$, \inserted{$B\equiv0\ (2)$}, dont on déduit
\[
\xi=\begin{pmatrix}A&B\\C&D\end{pmatrix},\qquad
\eta=\begin{pmatrix}R&S\\T&U\end{pmatrix},\qquad
\text{et }\nu\in\widehat{\mathbb Z}\text{ par}
\]
\[
\tag{157}
\begin{aligned}
C&=B-D+1,\qquad
A=\frac{1+BC}{D}=\frac1D(1+B^2-BD+B)\\
&&\inserted{{}=\mu(1+B+B^2)-B}
\end{aligned}
\]
\[
\tag{158}
\mu=\frac1D,\qquad\nu=\frac{\mu-1}{2},
\]
\deleted{\unsure{(159)}}
\[
\eta=\xi\ell_1^{-\nu}
=\begin{pmatrix}A&B\\C&D\end{pmatrix}
\begin{pmatrix}1&0\\-2\nu&1\end{pmatrix}
=\begin{pmatrix}A-2\nu B&B\\C-2\nu D&D\end{pmatrix}
\qquad\text{i.e.}
\]

\margin{on a normalisé le signe de $\eta$ comme \unsure{dit} dans (156)}

\[
\tag{159}
\text{i.e.}\quad
\unsure{U}=D,\quad S=\unsure{T}=B,\quad
R=A-2\nu B=\mu(1+B^2)=\frac{1+B^2}{D}
\]

Posons (comparer p.~\unsure{398'} (78))
\[
\tag{160}
\underline u_{\xi\deleted{\unsure{,\eta}}}
=\mu\begin{pmatrix}1&B\\0&D\end{pmatrix}
=\begin{pmatrix}\mu&\mu B\\0&1\end{pmatrix}
=\begin{pmatrix}\dfrac1D&B/D\\0&1\end{pmatrix}
\in\mathrm{Gl}(2,\widehat{\mathbb Z})
\]

\margin{\deleted{pour le moment on \unsure{dispose} d'un élément $\xi$ de $\unsure{\widetilde{\widehat{\Pi}}_{0,3}}$, \unsure{qui est déterminé|$\xi$ est déterminé par} comme élément de $\mathrm{Sl}(2,\widehat{\mathbb Z})'$ \unsure{ou dans|ou de} $\mathrm{Sl}(2,\widehat{\mathbb Z})$ \uuline{au signe près}}}

de sorte que la connaissance de $\underline u_\xi$ équivaut à celle \inserted{de $B$ et $D$, i.e. à celle} de $\xi\in\mathrm{Sl}(2,\widehat{\mathbb Z})$, ou encore à celle des couples $(\xi,\eta)$ ($\eta$ normalisé par $\eta=+\xi\ell_1^{-\nu}$). (Dans loc. cit. cet élément était noté $u_{\alpha,\beta}$... mais il ne dépend que de $\xi\in\widehat{\Pi}_{0,3}$, et $\overline{\mathrm m}$ seulement de $\underline\xi\in\mathrm{Sl}(2,\widehat{\mathbb Z})$ --- et sa connaissance équivaut à celle de $\underline\xi\in\mathrm{Sl}(2,\widehat{\mathbb Z})$ ---)

\margin{\uline{NB} On déduit $\xi,\eta$ de la connaissance de $\underline u$ par les formules
\deleted{\illegible{1}}
\[
\tag{161}
\left\{\begin{aligned}
\underline u\cdot{}^t\underline u&=\mu\eta\\
\mu&=\det\underline u
\end{aligned}\right.
\]
d'où
\[
\tag{162}
\left\{\begin{aligned}
\eta&=\mu^{-1}\underline u\cdot{}^t\underline u=\mu^{-1}\unsure{,}\\
\xi&=\eta\ell_1^{-\nu}=\mu^{-1}(\underline u\cdot{}^t\underline u)\ell_1^{-\nu}
\end{aligned}\right.
\]
}

\scanpage{222}

\deleted{\unsure{l'application}} On aura
\deleted{$\unsure{\mathcal M_{0,3}^{!\sim}(0)}$}
\deleted{\unsure{i.e. $\eta=(\underline u,\sigma_\infty)$}}
\[
(163)\qquad
\underline u_\xi(\sigma_\infty)
=\eta\sigma_\infty
=(\xi \ell_1^{-\nu})\sigma_\infty,
\qquad
\underline u_\xi(\rho)=\xi\rho
\]
\[
(164)\qquad
\underline u_\xi(\varepsilon_0)=\varepsilon_0^\mu
\qquad\text{et}\qquad
u_\xi(\ell_0]=\ell_0^\mu
\]

\margin{Les formules (163) s'écrivent aussi
\[
\unsure{(\underline u,\sigma_\infty)|[\underline u,\sigma_\infty]}=\eta,
\quad [\underline u,\rho]=\xi
\qquad (163\ \mathrm{bis}),
\]
elles donnent la façon de \unsure{reconstituer}
$(\unsure{\xi},\eta)$ à l'aide de $\underline u_\xi=\underline u$.}

Si on considère une matrice inversible quelconque
\[
\underline u=\begin{pmatrix}L&M\\N&P\end{pmatrix}
\in\mathrm{Gl}(2,\widehat{\mathbb Z})
\]
la relation
\[
\deleted{(\illegible{1})}\qquad
\underline u(\varepsilon_0)=\varepsilon_0^\mu
\quad\text{i.e.}\quad
u\varepsilon_0=\varepsilon_0^\mu
\qquad(\text{où }\mu\in\widehat{\mathbb Z}^{*})
\]
\deleted{\unsure{équivaut à}}\inserted{s'écrit}
\[
\begin{pmatrix}L&-L+M\\N&-N+P\end{pmatrix}
=\begin{pmatrix}L-\mu N&M-\mu P\\N&P\end{pmatrix}
\]
et équivaut à : \deleted{\unsure{$N$}}
\[
N=0,\qquad L=\mu P
\]
i.e.
\[
\underline u=\begin{pmatrix}\mu P&M\\0&P\end{pmatrix},
\qquad\text{d'où}\quad\det u=\mu P^2
\]

la condition supplémentaire
\[
\unsure{(165)}\qquad \det\underline u=\mu
\]
donne la condition
\[
\deleted{(\illegible{1})}\qquad P^2=\underline 1,
\]
la solution la plus \unsure{évidente} étant $P=1$,
ce qui donne \unsure{juste} les matrices
\[
\underline u=\begin{pmatrix}\mu&M\\0&1\end{pmatrix}
=\underline u_\xi
\]

\scanpage{223}
\mspage{411}

où $\xi=\begin{pmatrix}A&B\\C&D\end{pmatrix}$
\deleted{\unsure{satisfait}} est \unsure{donné|donnée} en termes de
$B,D$ par $D=\dfrac1\mu$, $B=DM=M/\mu$.

\uline{Théorème} L'application
\[
(165)\qquad
\begin{array}{rcl}
\mathcal M_{0,3}^{!\sim}[0]&\longrightarrow&
\mathrm{Gl}(2,\widehat{\mathbb Z})^{\deleted{\unsure{\prime|*}}}\\
\underset{u}{u_{\xi,\eta}}&\longmapsto&\underline u_\xi
\end{array},
\quad\text{ou encore }u\longmapsto\underline u
\]
$\underline u_\xi$ est donné par (160)) est un hom. de groupes.

\margin{\uline{NB} Du fait qu'on \unsure{s'est restreint à}
$\mathcal M_{0,3}^{!\sim}(0)$, on trouve un hom dans
$\mathrm{Gl}(2,\widehat{\mathbb Z})$,
pas seulement dans $\mathrm{Gl}(2,\widehat{\mathbb Z})/\mathord\pm1$.}

\margin{\uline{NB} Cette application se factorise par le
\unsure{ss-groupe|sous-groupe} triangulaire $\mathcal B$ ci-dessous (169)}

\uuline{Dém}
Considérons $u=u_{\xi,\eta}$,
\deleted{$u'=u_{\xi',\eta'}$, on aura \unsure{soit}
$u''=u_{\xi'',\eta''}$,}
$\underline u$ est son image dans $\mathrm{Gl}(2,\widehat{\mathbb Z})$,
et si $x\in\widehat\Pi_{0,3}$, $\underline x$
\deleted{$\xi''=u'(\xi)\xi'\qquad(\text{dans }\widehat\Pi_{0,3})$}
son image dans $\mathrm{Sl}(2,\widehat{\mathbb Z})$ :
\deleted{on aura \unsure{donc dans}}
\[
(166)\qquad
\underline{u(x)}=\underline u(\underline x)
\left(\stackrel{\text{déf}}{=}
\underline u\,\underline x\,\underline u^{-1}\right)
\]
car il suffit de le voir pour $x=\sigma_\infty$ et pour $x=\rho$,
où cette formule se réduit \unsure{à} (161). Ceci noté,
si $u=u_{\xi,\eta}\inserted{=u_\xi}$,
$u'=u_{\xi',\eta'}\inserted{=u_{\xi'}}$
sont deux éléments de $\mathcal M_{0,3}^{!\sim}(0)$,
$u''=u_{\xi'',\eta''}=u_{\xi''}$ leur produit, on aura
\inserted{(139)}
\[
\deleted{\unsure{(167)}}\qquad\xi''=u'(\xi)\xi'
\]
d'où
\[
\underline\xi''=\underline{u'(\xi)}\,\underline\xi^{\unsure{\prime}}
=\underline u'(\underline\xi)\,\underline\xi',\quad\text{i.e.}
\]
\[
\deleted{\unsure{(167)}}\quad(167)\qquad
\underline\xi''=\operatorname{int}
(\underline u_{\underline\xi'})(\underline\xi)\mathbin{.}\underline\xi'.
\]
et il suffit de vérifier que si trois solutions
\deleted{\illegible{1}}
\deleted{$\xi,\xi',\xi''$ des équations}
\[
\xi=\begin{pmatrix}A&B\\C&D\end{pmatrix},
\quad\xi'=\begin{pmatrix}A'&B'\\C'&D'\end{pmatrix},
\quad\xi''=\begin{pmatrix}A''&B''\\C''&D''\end{pmatrix}
\]
\unsure{des} équations
$AD-BC=1$, $C+D\unsure{-}B=1$, $D\in\widehat{\mathbb Z}^{*}$,
sont reliées par (167), alors on a
\[
(168)\qquad
\underline u_{\xi''}=\underline u_{\xi'}\mathbin{.}\underline u_\xi
\]

\scanpage{224}

\inserted{Comme $\xi$ et $\underline u=u_\xi$ se déterminent mutuellement
(163bis) il}
\deleted{\illegible{1}} revient au $\overline{\mathrm m}$ de dire que si,
sur le groupe
\[
(169)\qquad
\begin{array}{ll}
\mathcal B\subset\mathrm{Gl}(2,\widehat{\mathbb Z})&\\
\Vert&\\[-2pt]
\left\{\begin{pmatrix}\mu&M\\0&1\end{pmatrix}
\ \middle|\ \mu,M\in\widehat{\mathbb Z},\ \mu\text{ inv.}\right\}&
\end{array}
\]
on définit (cf (162))
\[
(170)\qquad
\left\{\begin{aligned}
F_\xi &: \mathcal B\longrightarrow\mathrm{Sl}(2,\widehat{\mathbb Z})\\
F_\xi(\underline u)
&=(\underbrace{\det\underline u}_{\mu})^{-1}
\underbrace{\deleted{\unsure{\underline u\,{}^t\underline u}}}
_{\inserted{[\underline u,\rho]}}
\overset{\deleted{\unsure{\mu-1=2\nu}}}
{\deleted{\unsure{\varepsilon_1^{-(\det\underline u-1)}|\ell_1^{-(\det\underline u-1)}}}}
\end{aligned}\right.
\]
alors on a la relation des 1-cocycles
\[
(171)\qquad
F_\xi(\underline u'\underline u)
=\operatorname{int}(\underline u')\bigl(F_\xi(\underline u)\bigr)
\mathbin{.}F_\xi(\underline u').
\]
Par multiplicativité de $\det$, on est \unsure{ramené}
\deleted{$\unsure{(\underline u,\underline u'\in\mathcal B)}$}
à voir que
\[
\underline u\longmapsto[\underline u,\rho]
=\underline u\,\rho(\underline u)^{-1}
\]
satisfait la relation des cocycles, ce qui est trivial.

\uline{Corollaire} L'application
\[
(172)\qquad
\begin{array}{ccl}
\mathcal M_{0,3}^{\unsure{!\sim}}
=\mathcal M_{0,3}^{!\sim}[0]\mathbin{.}
\widehat{\mathcal G}_{0,3}^{+}
&\longrightarrow&\mathrm{Gl}(2,\widehat{\mathbb Z})'\\
\deleted{\text{produit }\tfrac12\text{ direct}}&&\\
u_\xi\mathbin{.}a&\longrightarrow&(\underline u_\xi)'\mathbin{.}\underline a
\end{array}
\]
où $(\underline u_\xi)'$ \unsure{est la} réduction de $\underline u_\xi$
mod $\pm1$, et où $\underline a$ désigne l'image de $a$ dans
$\mathrm{Sl}(2,\widehat{\mathbb Z})'$) est un homomorphisme de groupes.

Résulte aussitôt du Th et des formules (163).

\margin{attention, ce n'est \uuline{pas} un produit semi-direct,
puisque $\mathcal M_{0,3}^{!\sim}(0)\cap
\widehat{\mathcal G}_{0,3}^{+}
=L_0\ne\{1\}$, il faut donc montrer que la formule de définition
ne dépend pas de la façon d'écrire un $u$ sous la forme
\unsure{$u_\xi a$} --- cela revient à voir que pour
$\lambda\in L_0$, \unsure{ces deux} images \illegible{2} \unsure{comme élément de l'un ou l'autre}
\unsure{facteur}, sont les $\overline{\mathrm m}$.}

\scanpage{225}
\mspage{412}

Dorénavant je vais \unsure{repérer|répérer} le couple
$(\underline\xi,\underline\eta)$ du point de vue adélique
$(\text{dans }\mathrm{Gl}(2,\widehat{\mathbb Z}))$ à l'aide
de $\underline u_\xi$, qui est donc de la forme
\[
(173)\qquad
\underline u=\begin{pmatrix}\mu&M\\0&1\end{pmatrix}
\]
\inserted{avec $M\equiv0\ (2)$ si $\underline u$ provient de
$\mathcal M_{0,3}^{!\sim}(0)$}
dont on tire $\xi,\eta$ par les formules \unsure{(163 bis)|(167 bis)},
\[
(174)\qquad
\underline\xi=[\underline u,\underline\rho],\qquad
\underline\eta=[\underline u,\underline\sigma_\infty].
\]

Pour $\underline u\in\mathcal M_{0,3}^{!\sim}(0)$, défini par
$(\xi,\eta)$, tant qu'on ne fait pas intervenir la
possibilité d'écrire $\xi$ et $\eta$ de façon unique comme
\[
\xi=\alpha\rho(\alpha)^{-1},\quad
\eta=\beta\sigma_\infty(\beta)^{-1}
\]
avec $\alpha\in\unsure{\widetilde{\widehat\Pi}_{0,3}|\widehat{\widetilde\Pi}_{0,3}}$,
$\beta\in\widehat\Pi_{0,3}$, les seuls invariants adéliques qu'on
puisse lui associer directement de façon naturelle est le couple
$(\underline\xi,\underline\eta)$, ou encore \unsure{(ce} qui revient au
$\overline{\mathrm m}$, mais est plus joli) $\underline u$.
Les éléments \deleted{\illegible{1}}
\[
\lambda=\ell_0^n\in L_0\hookrightarrow\mathcal M_{0,3}^{!\sim}[0]
\]
correspondent à
\[
\begin{aligned}
(175)\qquad
\xi=\eta=\ell_0^n\ell_1^{-n}
&=\begin{pmatrix}\unsure{1}&-2n\\0&\unsure{1}\end{pmatrix}
\begin{pmatrix}1&0\\-2n&1\end{pmatrix}\\
&=\begin{pmatrix}A=\unsure{1+4n^2}&B=-2n\\C=-2n&D=1\end{pmatrix}
\end{aligned}
\]
\margin{\unsure{prière} de faire la vérification \unsure{oubliée}
dont il \unsure{est} question dans la note \unsure{en marge}
\unsure{précédente}\ldots}
\[
(176)\qquad
\underline\lambda=\unsure{\underline{\underline{\ell_0}}^{\,n}|\underline{\ell_0}^{\,n}}
=\begin{pmatrix}1&-2n\\0&1\end{pmatrix}
\]
d'ailleurs on a
\[
(177)\qquad
\underbrace{\begin{pmatrix}1&-2n\\0&1\end{pmatrix}}_{\underline{\ell_0}^{\,n}}
\begin{pmatrix}\mu&M\\0&1\end{pmatrix}
=\begin{pmatrix}\mu&M-2n\\0&1\end{pmatrix}
\]
ce qui montre que \deleted{\illegible{1}}
\deleted{le quotient}\inserted{les quotients des}
sous-groupes $\mathcal B'_1$ de $\mathcal B$ (169) formés des
